\section{Planteamiento del Problema}

En la documentaci\'on de hardware t\'ecnico complejo, la comunicaci\'on visual sigue anclada con frecuencia en paradigmas est\'aticos, como archivos PDF, planos 2D y l\'aminas de ensamblaje, que disocian la informaci\'on visual de su contexto espacial. En esta propuesta, \textit{hardware complejo} se refiere a sistemas electromec\'anicos multicomponente con integraci\'on estructural, energ\'etica, electr\'onica y de control, alta densidad de relaciones espaciales entre piezas y necesidad de inspecci\'on visual detallada del ensamblaje. Esta desconexi\'on obliga a ingenieros y t\'ecnicos a realizar procesos mentales de reconstrucci\'on tridimensional que pueden incrementar la demanda cognitiva asociada a la tarea y consumir recursos de memoria de trabajo (Sweller, 1988).

Sin embargo, la migraci\'on hacia entornos 3D interactivos en la web enfrenta restricciones arquitect\'onicas significativas. Los modelos CAD de ingenier\'ia, basados en NURBS, ensamblajes extensos y activos de alta densidad geom\'etrica, no son directamente compatibles con las condiciones de despliegue de una aplicaci\'on Web orientada a acceso desde navegador. Aunque WebAssembly admite crecimiento din\'amico de memoria, la viabilidad real de Unity Web depende de controlar la huella inicial de los activos, la fragmentaci\'on del espacio de memoria del navegador y la necesidad de reservar bloques contiguos suficientemente grandes para heap y datos de aplicaci\'on (Unity Technologies, 2024a; WebAssembly Community Group, s. f.). En este contexto, persiste la necesidad de documentar y evaluar un \textit{pipeline} de \textit{Technical Art} que permita cerrar la brecha entre fidelidad visual, claridad t\'ecnica y rendimiento web sin sacrificar la legibilidad funcional del ensamblaje.

El problema t\'ecnico puede resumirse en la relaci\'on entre fidelidad y rendimiento:
\begin{itemize}
    \item \textbf{Huella inicial de datos.} Modelos CAD originales pueden alcanzar decenas o cientos de megabytes, y crecer a escalas mayores en ensamblajes extensos, por lo que requieren limpieza, simplificaci\'on geom\'etrica, \textit{baking} y empaquetado selectivo para ser viables en web.
    \item \textbf{Rendimiento en tiempo real.} Mantener 30 FPS o m\'as, equivalente a un \textit{frame time} de aproximadamente 33.33 ms, exige controlar presupuesto geom\'etrico, materiales, \textit{draw calls} y estados de render.
    \item \textbf{Latencia de interacci\'on.} Soluciones basadas en \textit{streaming} remoto pueden ofrecer alta fidelidad, pero introducen dependencia de red y latencias poco adecuadas para inspecci\'on t\'ecnica precisa.
\end{itemize}

Existe, por tanto, la necesidad de dise\~nar y evaluar un \textit{pipeline} de \textit{Technical Art} que permita desplegar modelos de hardware complejo en un visor 3D web con orientaci\'on espacial suficiente para tareas de inspecci\'on, costo computacional controlado y capacidad de an\'alisis visual del ensamblaje para prop\'ositos acad\'emicos e ingenieriles.

\subsection{Pregunta de investigaci\'on}

\noindent \textbf{Pregunta principal.} \textquestiondown Qu\'e diferencias descriptivas se observan en desempe\~no de tareas, usabilidad percibida y carga de trabajo percibida cuando usuarios con perfil af\'in realizan actividades de exploraci\'on e identificaci\'on del dron Holybro X500~V2 mediante un prototipo Web 3D interactivo, en comparaci\'on con un soporte 2D de referencia, y bajo qu\'e condiciones de rendimiento resulta t\'ecnicamente viable dicho prototipo en el navegador?

\newpage
